\documentclass{article}
\usepackage{fontspec} 
\usepackage{unicode-math}  % 先加载 unicode-math
\setmathfont{XITS Math}    % 再设置数学字体
%\usepackage[UTF8]{ctex} % 如果需要支持中文

\usepackage{anyfontsize}
\usepackage{xeCJK} % XeLaTeX 推荐 xeCJK，LuaLaTeX 也可用   
\usepackage{setspace}  % 添加这个包
\setstretch{1.2}      % 设置行距为1.5倍

\setCJKmainfont[
    Path = C:/USERS/11252/APPDATA/LOCAL/MICROSOFT/WINDOWS/FONTS/,  % 改用 Windows 字体目录
    UprightFont = {SourceHanSansCN-Light1.otf},  % 使用可变字体
    BoldFont = {SourceHanSansCN-Medium1.otf},
    LetterSpace = 2.0,  % 增加字间距
    WordSpace = 1.2,    % 增加词间距
    %BoldFeatures = {RawFeature={+wght=900}},
    %Scale=1
]{SourceHanSansCN}


\setmainfont[
    Path = C:/USERS/11252/APPDATA/LOCAL/MICROSOFT/WINDOWS/FONTS/,  % 改用 Windows 字体目录
    UprightFont = {SourceHanSansCN-Light1.otf},  % 使用可变字体
    BoldFont = {SourceHanSansCN-Medium1.otf},
    LetterSpace = 2.0,  % 增加字间距
    WordSpace = 1.2,    % 增加词间距
    %BoldFeatures = {RawFeature={+wght=900}},
    %Scale=1
]{SourceHanSansCN}

%\setCJKmainfont[
 %   Path = C:/USERS/11252/APPDATA/LOCAL/MICROSOFT/WINDOWS/FONTS/,
 %   UprightFont = {SourceHanSerifCN-Light-5.otf},
 %   BoldFont = {SourceHanSerifCN-Medium-6.otf},
%    Scale=1
%]{SourceHanSerifCN}

%\setmainfont{SourceHanSerif}  % 设置英文字体

% 处理冲突的命令
\let\oldbackepsilon\backepsilon
\let\backepsilon\relax
\let\oldeth\eth
\let\eth\relax
\let\olddigamma\digamma
\let\digamma\relax
\usepackage{float}

\usepackage[a4paper, left=0.1in, right=0.1in, top=0.05in, bottom=0.05in]{geometry} % 设置四边页边距为0.1英寸{geometry} % 设置页边距为0.5英寸
\usepackage{enumitem} % 用于自定义列表
\usepackage{amsmath}  % 数学公式支持
\usepackage{tikz}
\usetikzlibrary{arrows.meta, positioning}
\setlength{\parindent}{0pt} % 不缩进
\usepackage{etoolbox} % 用于列表操作
%调整类代码
\usepackage{comment}
\usepackage{tocloft} % 用于定制目录 样式
\usepackage{hyperref} %不需要目录盒子注释这
\usepackage{ifthen}
\usepackage{tikz}
\usetikzlibrary{arrows.meta}
\usepackage{titletoc}
\usepackage{graphicx}
\usepackage{float}

\usepackage{amssymb}  % 提供数学符号，包括\therefore
%目录设定
%快捷键 CMD + / 注释
%  \renewcommand{\cftdot}{} % 隐藏目录中的页码
%  \cftpagenumbersoff{section}
%  \cftpagenumbersoff{subsection}
%  \makeatletter
%  \renewcommand{\@pnumwidth}{0em}  % 设置页码宽度为0
%  \renewcommand{\@tocrmarg}{0em}   % 设置右边距为0
%  \makeatother
 




\usepackage{environ}

\newif\ifshowcontent  % 定义一个条件判断变量
\showcontenttrue    % 设置为 false，隐藏所有 question 环境的内容

\newcounter{question} % 题目计数器
\setcounter{question}{0}
\NewEnviron{question}{%
    \ifshowcontent
        \par\stepcounter{question}\textbf{\thequestion.} \BODY\par % 如果显示内容，则输出
    \fi
}

\NewEnviron{choices}{%
    \ifshowcontent
        \begin{enumerate}[label=\Alph*., leftmargin=*]
            \BODY
        \end{enumerate}\par
    \fi
}

\NewEnviron{correctanswer}{%
    \ifshowcontent
        \textbf{答案:}\BODY\par
    \fi
}



\newenvironment{explanation}
{\textbf{解析:}\par}
{\par}



% 新增考察方面命令
\newcommand{\checklist}[1]{
    \textbf{考察:} #1 \par % 在题目下方显示考察方面
    \addcontentsline{toc}{subsection}{题号\thequestion 考察: #1} % 将考察内容添加到目录
    \vspace{2em} % 添加1em的垂直空间
}

\graphicspath{{images/}} %统一


\begin{document}
 %\fontsize{22}{31}\selectfont
 \tableofcontents % 添加目录
\newpage
%\pagestyle{empty}




\section{考点1:全球金融危机前的资本监管}
\setcounter{question}{0}


\begin{question}
    As a result of the new Basel standards, every bank must now calculate explicit capital charges to cover operational risk using one of three approaches: the basic indicator approach (BIA), the standardized approach (SA), and the advanced measurement approach (AMA). Which of the following statements are true with respect to these operational risk approaches?
    I. In practice the AMA is the most stringent approach for operational risk
    II. The most popular method to satisfy the AMA is the loss distribution approach
    III. The AMA allows a bank to build its own operational risk model and measurement system comparable to market risk standards
    IV. BIA is widely used in insurance and actuarial science
\end{question}

\begin{choices}
    \item II only
    \item III and IV
    \item II and III
    \item None are correct
\end{choices}

\begin{correctanswer}
    C
\end{correctanswer}

\begin{explanation}
    \textbf{A选项错误。}
    这个说法不完整：
    \begin{itemize}
        \item II项正确：LDA是满足AMA最常用的方法
        \item 但遗漏了III项正确说法
    \end{itemize}

    \textbf{B选项错误。}
    这个说法不准确：
    \begin{itemize}
        \item III项正确：AMA允许银行建立自己的风险模型
        \item IV项错误：LDA而非BIA用于保险和精算
    \end{itemize}

    \textbf{C选项正确。}
    这个说法正确：
    \begin{itemize}
        \item II项正确：LDA是满足AMA最常用的方法
        \item III项正确：AMA允许银行建立自己的风险模型
        \item I项错误：AMA是最灵活的方法
        \item IV项错误：LDA而非BIA用于保险和精算
    \end{itemize}

    \textbf{D选项错误。}
    这个说法不准确：
    \begin{itemize}
        \item II和III项都是正确的
        \item 不应选择"None are correct"
    \end{itemize}
\end{explanation}

\checklist{巴塞尔操作风险计量方法的特点}


\begin{question}
    According to the Basel Committee, which of the options below is not a quantitative standard that a bank must meet before it is permitted to use the Advanced Measurement Approach (AMA) for operational risk capital?
\end{question}

\begin{choices}
    \item A bank's risk measurement system should be sufficiently 'granular' to capture the major drivers of operational risk affecting the shape of the tail of the loss estimates
    \item Supervisors will require the bank to calculate its regulatory capital as the Unexpected Loss (UL) disregarding the Expected Losses (EL)
    \item Internally generated operational risk measures used for regulatory capital purposes must be based on a minimum 5-year observation period of loss data. When the bank first moves to the AMA, a 3-year historical data window is acceptable
    \item The tracking of internal loss data
\end{choices}

\begin{correctanswer}
    B
\end{correctanswer}

\begin{explanation}
    \textbf{A选项错误。}
    这个说法不准确：
    \begin{itemize}
        \item 风险计量系统需要足够细化
        \item 这是AMA的定量标准之一
    \end{itemize}

    \textbf{B选项正确。}
    这个说法正确：
    \begin{itemize}
        \item 监管资本要求应包括EL和UL
        \item 除非银行能证明已充分覆盖EL
        \item 这不是定量标准
    \end{itemize}

    \textbf{C选项错误。}
    这个说法不准确：
    \begin{itemize}
        \item 需要至少5年损失数据
        \item 首次采用时可接受3年数据
        \item 这是定量标准之一
    \end{itemize}

    \textbf{D选项错误。}
    这个说法不准确：
    \begin{itemize}
        \item 内部损失数据跟踪是要求
        \item 这是定量标准之一
    \end{itemize}
\end{explanation}

\checklist{AMA操作风险计量的定量标准}


\begin{question}
    Compare the capital requirements of Bank A and Bank B based on the following equal-sized portfolios.

    \begin{center}
    \begin{tabular}{l@{\hspace{2cm}}l}
        \textbf{Bank A} & \textbf{Bank B} \\
        30\% OECD sovereign debt & 60\% OECD sovereign debt \\
        70\% non-OECD sovereign debt & 40\% non-OECD sovereign debt \\
    \end{tabular}
    \end{center}
    
    The credit risk capital requirement of Bank A will be greater than the credit risk capital requirement of Bank B under the:
\end{question}

\begin{choices}
    \item standardized approach of Basel I
    \item standardized approach of Basel II
    \item IRB foundation approach of Basel II
    \item advanced measurement approach of Basel II
\end{choices}

\begin{correctanswer}
    A
\end{correctanswer}

\begin{explanation}
    \textbf{A选项正确。}
    这个说法正确：
    \begin{itemize}
        \item 巴塞尔I中OECD国家风险权重低于非OECD国家
        \item A银行非OECD债务比例更高
        \item 因此A银行资本要求更高
    \end{itemize}

    \textbf{B选项错误。}
    这个说法不准确：
    \begin{itemize}
        \item 巴塞尔II不再简单区分OECD和非OECD
        \item 基于外部评级确定风险权重
    \end{itemize}

    \textbf{C选项错误。}
    这个说法不准确：
    \begin{itemize}
        \item 基础IRB法基于内部评级
        \item 需要更多信息才能比较
    \end{itemize}

    \textbf{D选项错误。}
    这个说法不准确：
    \begin{itemize}
        \item 高级法完全基于内部模型
        \item 需要更多信息才能比较
    \end{itemize}
\end{explanation}

\checklist{巴塞尔协议下资本要求的比较}


\begin{question}
    Michael Chen, FRM, is a consultant for U.S.-based Metropolitan Bank. Chen attended a meeting where a Senior Vice President made the following statements about the Basel II Accord:
    I. By switching from the standardized approach to the foundation IRB approach, our risk weightings for a majority of the bank's assets are lower, which could reduce our capital requirements by as much as 20% next year
    II. Under the IRB advanced approach, we generate all the estimates used in the models
    III. Pillar 2 concerns external monitoring and supervisory review
    How many of the statements are correct?
\end{question}

\begin{choices}
    \item None
    \item One
    \item Two
    \item Three
\end{choices}

\begin{correctanswer}
    B
\end{correctanswer}

\begin{explanation}
    \textbf{A选项错误。}
    这个说法不准确：
    \begin{itemize}
        \item II项正确
        \item 高级IRB法确实由银行自行估计所有参数
    \end{itemize}

    \textbf{B选项正确。}
    这个说法正确：
    \begin{itemize}
        \item I项错误：IRB法有过渡期限制
        \item 第一年不低于前一年的90%
        \item 第二年不低于前一年的80%
        \item II项正确：高级IRB法由银行自行估计
        \item III项错误：第二支柱是监管审查
        \item 外部监督是第三支柱的内容
    \end{itemize}

    \textbf{C选项错误。}
    这个说法不准确：
    \begin{itemize}
        \item 只有II项正确
        \item I和III项都错误
    \end{itemize}

    \textbf{D选项错误。}
    这个说法不准确：
    \begin{itemize}
        \item 只有II项正确
        \item I和III项都错误
    \end{itemize}
\end{explanation}

\checklist{巴塞尔II协议的关键特征}


\begin{question}
    As a risk manager of Bank XYZ, Michael is asked to calculate the market risk capital charge of the bank's trading portfolio under the internal models approach using the information given in the table below. Assuming the return of the bank's trading portfolio is normally distributed, what is the market risk capital charge of the trading portfolio?

    \begin{table}[htbp]
        \centering
        \begin{tabular}{l r}
            \hline
            VaR (95\%, 1-day) of last trading day                & USD 45,000 \\
            Average VaR (95\%, 1-day) for last 60 trading days   & USD 30,000 \\
            Multiplication Factor                                & 2.5 \\
            \hline
        \end{tabular}
        \caption{VaR 相关数据}
    \end{table}
\end{question}

\begin{choices}
    \item USD 95,000
    \item USD 145,000
    \item USD 195,000
    \item USD 250,000
\end{choices}

\begin{correctanswer}
    D
\end{correctanswer}

\begin{explanation}
    \textbf{A选项错误。}
    这个说法不准确：
    \begin{itemize}
        \item 未考虑时间平方根法则
        \item 未考虑置信水平转换
    \end{itemize}

    \textbf{B选项错误。}
    这个说法不准确：
    \begin{itemize}
        \item 计算过程有误
        \item 未正确应用乘数因子
    \end{itemize}

    \textbf{C选项错误。}
    这个说法不准确：
    \begin{itemize}
        \item 计算过程有误
        \item 未正确考虑最大值
    \end{itemize}

    \textbf{D选项正确。}
    这个说法正确：
    \begin{itemize}
        \item 市场风险资本要求 = MAX(45,000 × √10/1.65 × 2.326, 2.5 × 30,000 × √10/1.65 × 2.326)
        \item 需要将1天95\% VaR转换为10天95\% VaR
        \item 考虑乘数因子和最大值
    \end{itemize}
\end{explanation}

\checklist{市场风险资本要求的计算方法}


\begin{question}
    Basel II/III contains capital rules for the banking industry, while Solvency II looks at the regulatory framework and quantifiable risks in the insurance industry. Comparing Solvency II and Basel II/III, which statement best characterizes the Basel II/III approach?
\end{question}

\begin{choices}
    \item Basel II/III does not try to achieve safety for the entire company, but rather focuses on asset-specific risks
    \item Basel II/III requires a more holistic perspective than Solvency II
    \item Basel II/III does not take in consideration pro-cyclical effects
    \item The goal of Basel II/III is to operate under no supervisory authority
\end{choices}

\begin{correctanswer}
    A
\end{correctanswer}

\begin{explanation}
    \textbf{A选项正确。}
    这个说法正确：
    \begin{itemize}
        \item 巴塞尔II/III更关注系统性风险
        \item 专注于三大资产特定风险类别
        \item 市场风险、信用风险和操作风险
    \end{itemize}

    \textbf{B选项错误。}
    这个说法不准确：
    \begin{itemize}
        \item Solvency II更注重整体视角
        \item 巴塞尔II/III更关注具体风险
    \end{itemize}

    \textbf{C选项错误。}
    这个说法不准确：
    \begin{itemize}
        \item 巴塞尔II/III考虑了顺周期性
        \item 包含逆周期资本缓冲要求
    \end{itemize}

    \textbf{D选项错误。}
    这个说法不准确：
    \begin{itemize}
        \item 巴塞尔II/III强调监管审查
        \item 第二支柱专门关注监管
    \end{itemize}
\end{explanation}

\checklist{巴塞尔II/III与Solvency II的比较}


\begin{question}
    A senior market risk analyst at bank ABC is examining Basel Committee's Fundamental Review of the Trading Book (FRTB) rules for calculating capital requirement. ABC would like to use an internal models approach to calculating market risk capital, so that analyst focuses on the FRTB guidelines for calculating ES to determine market risk capital using this approach. Which of the following statements correctly describes FRTB guidance when implementing an internal models approach?
\end{question}

\begin{choices}
    \item If a bank determines that some risk factors cannot be modeled, the required capital will be equal to 1.5 times the prior day's ES
    \item A bank can choose the most advantageous way to estimate ES, either at the portfolio level or at the trading desk level
    \item If historical data is not available on certain risk factors, a bank can estimate ES using an approved process of scaling up ES estimated with available data
    \item A bank can estimate ES by using a cascade approach that averages ten rolling ES estimates during a historically stressed period
\end{choices}

\begin{correctanswer}
    C
\end{correctanswer}

\begin{explanation}
    \textbf{A选项错误。}
    这个说法不准确：
    \begin{itemize}
        \item 无法建模的风险因素应使用标准法
        \item 不是简单乘以1.5倍
    \end{itemize}

    \textbf{B选项错误。}
    这个说法不准确：
    \begin{itemize}
        \item 银行不能自由选择估计方法
        \item 必须遵循FRTB的具体要求
    \end{itemize}

    \textbf{C选项正确。}
    这个说法正确：
    \begin{itemize}
        \item 允许使用经批准的放大方法
        \item 基于可用数据估计ES
        \item 符合FRTB的灵活性要求
    \end{itemize}

    \textbf{D选项错误。}
    这个说法不准确：
    \begin{itemize}
        \item 级联方法不是FRTB推荐的方法
        \item 未提及10个滚动估计
    \end{itemize}
\end{explanation}

\checklist{FRTB内部模型法的实施要求}




\newpage



\section{考点2:全球金融危机后的偿付能力、流动性和其他监管问题}
\setcounter{question}{0}
\begin{question}
    In the latest guidelines for computing capital for incremental risk in the trading book, the incremental risk charge (IRC) addresses a number of perceived shortcomings in the 99%/10-day VaR framework. Which of the following statements about the IRC are correct?
    I. For all IRC-covered positions, the IRC model must measure losses due to default and migration over a one-year horizon at a 99% confidence level
    II. A bank can incorporate into its IRC model any securitization positions that hedge underlying credit instruments held in the trading account
    III. A bank must calculate the IRC measure at least weekly, or more frequently as directed by its supervisor
    IV. The incremental risk capital charge is the maximum of (1) the average of the IRC measures over 12 weeks and (2) the most recent IRC measure
\end{question}

\begin{choices}
    \item I and II
    \item III and IV
    \item I, II, and III
    \item II, III, and IV
\end{choices}

\begin{correctanswer}
    B
\end{correctanswer}

\begin{explanation}
    \textbf{A选项错误。}
    这个说法不准确：
    \begin{itemize}
        \item I项错误：置信水平应为99.9%
        \item II项错误：证券化头寸适用银行账户资本要求
    \end{itemize}

    \textbf{B选项正确。}
    这个说法正确：
    \begin{itemize}
        \item III项正确：至少每周计算IRC
        \item IV项正确：取12周平均值和最新值的最大值
        \item I项错误：置信水平应为99.9%
        \item II项错误：证券化头寸不适用IRC
    \end{itemize}

    \textbf{C选项错误。}
    这个说法不准确：
    \begin{itemize}
        \item I和II项都错误
        \item III项正确
    \end{itemize}

    \textbf{D选项错误。}
    这个说法不准确：
    \begin{itemize}
        \item II项错误
        \item III和IV项正确
    \end{itemize}
\end{explanation}

\checklist{增量风险资本要求(IRC)的关键特征}


\begin{question}
    In March 2009, the Basel Committee published the consultative document "Guidelines for computing capital for Incremental risk in the trading book". The incremental risk charge (IRC) defined in that document aims to complement additional standards being applied to the value-at-risk modeling framework which address a number of perceived shortcomings in the 99%/10-day VaR framework. Which of the following statements about the IRC is/are correct?
    I. For all IRC-covered positions, a bank's IRC model must measure losses due to default and migration over a one-year capital horizon at a 99% confidence level
    II. A bank must calculate the IRC measure at least weekly, or more frequently as directed by its supervisor
\end{question}

\begin{choices}
    \item Statement I only
    \item Statement II only
    \item Both statements are correct
    \item Both statements are incorrect
\end{choices}

\begin{correctanswer}
    B
\end{correctanswer}

\begin{explanation}
    \textbf{A选项错误。}
    这个说法不准确：
    \begin{itemize}
        \item I项错误：置信水平应为99.9%
        \item 需要考虑流动性期限
    \end{itemize}

    \textbf{B选项正确。}
    这个说法正确：
    \begin{itemize}
        \item I项错误：置信水平应为99.9%
        \item 需要考虑流动性期限
        \item II项正确：至少每周计算IRC
        \item 可根据监管要求更频繁计算
    \end{itemize}

    \textbf{C选项错误。}
    这个说法不准确：
    \begin{itemize}
        \item I项错误
        \item II项正确
    \end{itemize}

    \textbf{D选项错误。}
    这个说法不准确：
    \begin{itemize}
        \item I项错误
        \item II项正确
    \end{itemize}
\end{explanation}

\checklist{增量风险资本要求(IRC)的基本要求}


\begin{question}
    In recent years, large dealer banks financed significant fractions of their assets using short-term, often overnight, repurchase (repo) agreements in which creditors held bank securities as collateral against default losses. The table below shows the quarter-end financing of four broker-dealer banks. All values are in USD billions:

    \begin{table}[htbp]
        \centering
        \begin{tabular}{lcccc}
            \hline
            & \textbf{Bank A} & \textbf{Bank B} & \textbf{Bank C} & \textbf{Bank D} \\
            \hline
            Financial instruments owned & 823 & 629 & 723 & 382 \\
            Pledged as collateral       & 272 & 289 & 380 & 155 \\
            \hline
        \end{tabular}
        \caption{各银行金融工具持有及质押情况}
    \end{table}
    In the event that repo creditors become nervous about a bank's solvency, which bank is least vulnerable to a liquidity crisis?
\end{question}

\begin{choices}
    \item Bank A
    \item Bank B
    \item Bank C
    \item Bank D
\end{choices}

\begin{correctanswer}
    A
\end{correctanswer}

\begin{explanation}
    \textbf{A选项正确。}
    这个说法正确：
    \begin{itemize}
        \item 质押比例最低：272/823 = 33%
        \item 对短期回购融资依赖最小
        \item 流动性风险最低
    \end{itemize}

    \textbf{B选项错误。}
    这个说法不准确：
    \begin{itemize}
        \item 质押比例：289/629 = 46%
        \item 高于A银行
    \end{itemize}

    \textbf{C选项错误。}
    这个说法不准确：
    \begin{itemize}
        \item 质押比例：380/723 = 53%
        \item 高于A银行
    \end{itemize}

    \textbf{D选项错误。}
    这个说法不准确：
    \begin{itemize}
        \item 质押比例：155/382 = 41%
        \item 高于A银行
    \end{itemize}
\end{explanation}

\checklist{银行流动性风险分析}


\begin{question}
    The balance sheet for Pierce Bankholdings as of December 31, 2013, included the following items:
    \begin{itemize}
        \item Preferred stock (cumulative): \$200,000
        \item Common stock: \$1,800,000
        \item Retained earnings: \$16,000,000
        \item Unrealized gains on long-term equity holdings: \$75,000
    \end{itemize}
    Based only on this information, estimate the Tier 1 and Tier 2 capital of Pierce Bankholdings as of 12/31/13.
\end{question}

\begin{choices}
    \item \$1,800,000 and \$200,000
    \item \$1,800,000 and \$275,000
    \item \$17,800,000 and \$200,000
    \item \$17,800,000 and \$275,000
\end{choices}

\begin{correctanswer}
    D
\end{correctanswer}

\begin{explanation}
    \textbf{A选项错误。}
    这个说法不准确：
    \begin{itemize}
        \item 未计入留存收益
        \item 未计入未实现收益
    \end{itemize}

    \textbf{B选项错误。}
    这个说法不准确：
    \begin{itemize}
        \item 未计入留存收益
        \item 未实现收益计入正确
    \end{itemize}

    \textbf{C选项错误。}
    这个说法不准确：
    \begin{itemize}
        \item 一级资本计算正确
        \item 未计入未实现收益
    \end{itemize}

    \textbf{D选项正确。}
    这个说法正确：
    \begin{itemize}
        \item 一级资本 = 普通股 + 留存收益 = \$17,800,000
        \item 二级资本 = 累积优先股 + 未实现收益 = \$275,000
    \end{itemize}
\end{explanation}

\checklist{银行资本结构分析}


\begin{question}
    \textbf{Large Bank} is attempting to transition to the new Basel III standards. Specifically, they are wondering if their liquidity and funding ratios meet the updated requirements as specified by the Basel Committee. Given the following information, what is the bank's current liquidity coverage ratio?

    \begin{itemize}
        \item High-quality liquid assets: \$300
        \item Marketable securities: \$125
        \item Required amount of stable funding: \$250
        \item Cash inflows over the next 30 days: \$214
        \item Net cash outflows over the next 30 days: \$285
        \item Long-term economic capital: \$500
        \item Available amount of stable funding: \$255
    \end{itemize}
\end{question}

\begin{choices}
    \item 95\%
    \item 105\%
    \item 125\%
    \item 140\%
\end{choices}

\begin{correctanswer}
    B
\end{correctanswer}

\begin{explanation}
    \textbf{A选项错误。}
    这个说法不准确：
    \begin{itemize}
        \item 计算错误
        \item 低于巴塞尔III要求
    \end{itemize}

    \textbf{B选项正确。}
    这个说法正确：
    \begin{itemize}
        \item LCR = 高质量流动资产/30天净现金流出
        \item LCR = \$300/\$285 = 105.3\%
        \item 满足巴塞尔III最低100\%要求
    \end{itemize}

    \textbf{C选项错误。}
    这个说法不准确：
    \begin{itemize}
        \item 计算错误
        \item 高估了LCR
    \end{itemize}

    \textbf{D选项错误。}
    这个说法不准确：
    \begin{itemize}
        \item 计算错误
        \item 高估了LCR
    \end{itemize}
\end{explanation}

\checklist{流动性覆盖率(LCR)的计算}


\begin{question}
    A common trade during 2004 and 2005 was to sell protection on the equity tranche and buy protection of the mezzanine tranche of the CDX.NA.IG index. Which of the following statements regarding this trade is least accurate?
\end{question}

\begin{choices}
    \item The trade was set up to be default-risk neutral at initiation
    \item The trade was short credit spread risk on the equity tranche and long credit spread risk on the mezzanine tranche
    \item The main motivation for the trade was to achieve a positively convex payoff profile
    \item The trade was designed to benefit from credit spread volatilities
\end{choices}

\begin{correctanswer}
    B
\end{correctanswer}

\begin{explanation}
    \textbf{A选项错误。}
    这个说法不准确：
    \begin{itemize}
        \item 该说法正确
        \item 交易在开始时是违约风险中性的
    \end{itemize}

    \textbf{B选项正确。}
    这个说法正确：
    \begin{itemize}
        \item 该说法错误
        \item 实际是：对权益层做多信用和信用利差风险
        \item 对夹层做空信用和信用利差风险
        \item 方向完全相反
    \end{itemize}

    \textbf{C选项错误。}
    这个说法不准确：
    \begin{itemize}
        \item 该说法正确
        \item 主要目的是获得正凸性收益
    \end{itemize}

    \textbf{D选项错误。}
    这个说法不准确：
    \begin{itemize}
        \item 该说法正确
        \item 交易设计是为了从信用利差波动中获利
    \end{itemize}
\end{explanation}

\checklist{CDO交易策略分析}


\begin{question}
    The OTC derivatives market was considered by many to have been partly responsible for the 2008 credit crisis. When the G20 leaders met in Pittsburgh in September 2009 in the aftermath of the 2008 crisis, they wanted to reduce systemic risk by regulating the OTC market. Which of the following is not a post-crisis change in regulation of OTC markets?
\end{question}

\begin{choices}
    \item All standardized OTC derivatives be cleared through CCPs
    \item Standardized OTC derivatives be traded on electronic platforms
    \item All trades in the OTC market be reported to a central trade repository
    \item OTC derivatives can be settled separately
\end{choices}

\begin{correctanswer}
    D
\end{correctanswer}

\begin{explanation}
    \textbf{A选项错误。}
    这个说法不准确：
    \begin{itemize}
        \item 该说法正确
        \item 要求所有标准化OTC衍生品通过CCP清算
        \item 目的是降低系统性风险
    \end{itemize}

    \textbf{B选项错误。}
    这个说法不准确：
    \begin{itemize}
        \item 该说法正确
        \item 要求标准化OTC衍生品在电子平台交易
        \item 目的是提高透明度
    \end{itemize}

    \textbf{C选项错误。}
    这个说法不准确：
    \begin{itemize}
        \item 该说法正确
        \item 要求所有OTC市场交易报告给中央交易库
        \item 为监管提供重要信息
    \end{itemize}

    \textbf{D选项正确。}
    这个说法正确：
    \begin{itemize}
        \item 该说法错误
        \item 危机后改革要求集中清算
        \item 不允许单独结算
    \end{itemize}
\end{explanation}

\checklist{OTC衍生品市场监管改革}




\newpage


\section{考点3:巴塞尔II改革摘要}
\setcounter{question}{0}
\begin{question}
    The CFO at a bank is preparing a report to the board of directors on its compliance with Basel requirements. The bank's average capital and total exposure for the most recent quarter is as follows:

    \begin{table}[htbp]
        \centering
        \begin{tabular}{|l|r|}
            \hline
            \textbf{REGULATORY CAPITAL} & \textbf{USD MILLIONS} \\
            \hline
            Total Common Equity Tier 1 Capital & 108 \\
            \hline
            Additional Tier 1 Capital & 28 \\
            \quad Prior to regulatory adjustments & 34 \\
            \quad Regulatory adjustments & 6 \\
            \hline
            Total Tier 1 Capital & 136 \\
            \hline
            Tier 2 Capital & 36 \\
            \quad Prior to regulatory adjustments & 45 \\
            \quad Regulatory adjustments & 9 \\
            \hline
            Total Capital & 172 \\
            Total Average Exposure & 3678 \\
            \hline
        \end{tabular}
        \caption{监管资本结构表}
    \end{table}
    Using the Basel III framework, which of the following is the best estimate of the bank's current leverage ratio?
\end{question}

\begin{choices}
    \item 2.94\%
    \item 3.70\%
    \item 4.68\%
    \item 5.08\%
\end{choices}

\begin{correctanswer}
    B
\end{correctanswer}

\begin{explanation}
    \textbf{A选项错误。}
    这个说法不准确：
    \begin{itemize}
        \item 计算错误
        \item 使用了错误的资本基础
    \end{itemize}

    \textbf{B选项正确。}
    这个说法正确：
    \begin{itemize}
        \item 杠杆率 = 一级资本/总敞口
        \item 杠杆率 = 136/3,678 = 3.70\%
        \item 符合巴塞尔III要求
    \end{itemize}

    \textbf{C选项错误。}
    这个说法不准确：
    \begin{itemize}
        \item 计算错误
        \item 使用了总资本而非一级资本
    \end{itemize}

    \textbf{D选项错误。}
    这个说法不准确：
    \begin{itemize}
        \item 计算错误
        \item 使用了错误的资本基础
    \end{itemize}
\end{explanation}

\checklist{银行杠杆率计算}



\begin{question}
    The CFO at a bank is preparing a report to the board of directors on its compliance with Basel requirements. The bank's average capital and total exposure for the most recent quarter is as follows:

    \begin{table}[htbp]
        \centering
        \begin{tabular}{|l|r|}
            \hline
            \textbf{Regulatory Capital} & \textbf{USD Millions} \\
            \hline
            Total Common Equity Tier 1 Capital & 108 \\
            Additional Tier 1 Capital & 28 \\
            \quad Prior to regulatory adjustments & 34 \\
            \quad Regulatory Adjustments & 6 \\
            Total Tier 1 Capital & 136 \\
            Tier 2 capital & 36 \\
            \quad Prior to regulatory adjustments & 45 \\
            \quad Regulatory adjustments & 9 \\
            Total Capital & 172 \\
            Total Average Exposure & 3678 \\
            \hline
        \end{tabular}
        \caption{监管资本结构表}
    \end{table}

    Using the Basel III framework, which of the following is the best estimate of the bank's current leverage ratio?
\end{question}

\begin{choices}
    \item 2.94\%
    \item 3.70\%
    \item 4.68\%
    \item 5.08\%
\end{choices}

\begin{correctanswer}
    B
\end{correctanswer}

\begin{explanation}
    \textbf{A选项错误。}
    这个说法不准确：
    \begin{itemize}
        \item 计算错误
        \item 使用了错误的资本基础
    \end{itemize}

    \textbf{B选项正确。}
    这个说法正确：
    \begin{itemize}
        \item 杠杆率 = 一级资本/总敞口
        \item 杠杆率 = 136/3,678 = 3.70\%
        \item 符合巴塞尔III要求
    \end{itemize}

    \textbf{C选项错误。}
    这个说法不准确：
    \begin{itemize}
        \item 计算错误
        \item 使用了总资本而非一级资本
    \end{itemize}

    \textbf{D选项错误。}
    这个说法不准确：
    \begin{itemize}
        \item 计算错误
        \item 使用了错误的资本基础
    \end{itemize}
\end{explanation}

\checklist{银行杠杆率计算}

\begin{question}
    The Basel III reforms introduced a leverage buffer ratio for all:
\end{question}

\begin{choices}
    \item global banks
    \item unregulated global banks
    \item global systemically important banks
    \item banks regulated by the U.S. Federal Reserve and the European Banking Authority
\end{choices}

\begin{correctanswer}
    C
\end{correctanswer}

\begin{explanation}
    \textbf{A选项错误。}
    这个说法不准确：
    \begin{itemize}
        \item 杠杆缓冲要求不适用于所有全球银行
        \item 仅适用于G-SIBs
    \end{itemize}

    \textbf{B选项错误。}
    这个说法不准确：
    \begin{itemize}
        \item 未受监管的全球银行不适用
        \item 杠杆缓冲要求针对受监管的G-SIBs
    \end{itemize}

    \textbf{C选项正确。}
    这个说法正确：
    \begin{itemize}
        \item 巴塞尔III改革为G-SIBs引入杠杆缓冲要求
        \item 必须由一级资本满足
        \item 额外要求1\%-3.5\%的一级核心资本
        \item 用于吸收潜在损失
    \end{itemize}

    \textbf{D选项错误。}
    这个说法不准确：
    \begin{itemize}
        \item 不仅限于美欧监管的银行
        \item 适用于所有G-SIBs
    \end{itemize}
\end{explanation}

\checklist{巴塞尔III杠杆缓冲要求}




\newpage



\section{考点4: 巴塞尔III：完成危机后改革}
\setcounter{question}{0}
\begin{question}
    All the following are operational risk loss events, except:
\end{question}

\begin{choices}
    \item An individual shows up at a branch presenting a check written by a customer for an amount substantially exceeding the customer's low checking account balance. When the bank calls the customer to ask him for the funds, the phone is disconnected and the bank cannot recover the funds
    \item A bank, acting as a trustee for a loan pool, receives less than the projected funds due to delayed repayment of certain loans
    \item During an adverse market movement, the computer network system becomes overwhelmed, and only intermittent pricing information is available to the bank's trading desk, leading to large losses as traders become unable to alter their hedges in response to falling prices
    \item A loan officer inaccurately enters client financial information into the bank's proprietary credit risk model
\end{choices}

\begin{correctanswer}
    B
\end{correctanswer}

\begin{explanation}
    \textbf{A选项错误。}
    这个说法不准确：
    \begin{itemize}
        \item 该事件属于操作风险
        \item 属于外部欺诈类型
    \end{itemize}

    \textbf{B选项正确。}
    这个说法正确：
    \begin{itemize}
        \item 该事件不属于操作风险
        \item 属于信用风险
        \item 是由于贷款延迟还款导致的损失
    \end{itemize}

    \textbf{C选项错误。}
    这个说法不准确：
    \begin{itemize}
        \item 该事件属于操作风险
        \item 属于系统故障类型
    \end{itemize}

    \textbf{D选项错误。}
    这个说法不准确：
    \begin{itemize}
        \item 该事件属于操作风险
        \item 属于内部流程失败类型
    \end{itemize}
\end{explanation}

\checklist{操作风险事件识别}


\begin{question}
    The Basel Committee requires that four data elements be used to calculate a bank's operational risk capital charge. Which of the following is not one of the four elements?
\end{question}

\begin{choices}
    \item External data
    \item Scenario analysis
    \item Internal loss data
    \item Regression analysis and other statistical tools
\end{choices}

\begin{correctanswer}
    D
\end{correctanswer}

\begin{explanation}
    \textbf{A选项错误。}
    这个说法不准确：
    \begin{itemize}
        \item 外部数据是四个要素之一
        \item 用于补充内部数据不足
    \end{itemize}

    \textbf{B选项错误。}
    这个说法不准确：
    \begin{itemize}
        \item 情景分析是四个要素之一
        \item 用于评估潜在风险
    \end{itemize}

    \textbf{C选项错误。}
    这个说法不准确：
    \begin{itemize}
        \item 内部损失数据是四个要素之一
        \item 反映银行自身历史损失
    \end{itemize}

    \textbf{D选项正确。}
    这个说法正确：
    \begin{itemize}
        \item 回归分析不是四个要素之一
        \item 四个要素是：
        \item 1. 内部损失数据
        \item 2. 外部数据
        \item 3. 情景分析
        \item 4. 业务环境和内部控制因素(BEICFs)
    \end{itemize}
\end{explanation}

\checklist{操作风险资本要求的计算要素}


\begin{question}
    A senior compliance officer at a multinational bank is presenting to a group of analysts on the Basel Committee's approaches for calculating market risk capital. The compliance officer describes the evolution of the Basel framework for market risk capital calculations and focuses on the Fundamental Review of the Trading Book(FRTB). Which of the following statements correctly describes an approach feature of the FRTB?
\end{question}

\begin{choices}
    \item Permission to use an internal models approach is granted on a bank-wide basis, however, each trading desk is responsible for carrying out its own market risk capital calculations
    \item Banks can determine the most appropriate liquidity horizons to use for market variables in market risk capital calculations, rather than relying on guidance in the Basel framework
    \item Compared to previous Basel frameworks, banks are given more flexibility in allocating instruments to either the trading book or the banking book, potentially reducing the volatility of their calculated market risk capital
    \item The standardized approach for calculating market risk capital must be used to provide a floor for capital requirements, even if a bank is approved to use an internal models approach
\end{choices}

\begin{correctanswer}
    D
\end{correctanswer}

\begin{explanation}
    \textbf{A选项错误。}
    这个说法不准确：
    \begin{itemize}
        \item 内部模型法按交易台审批
        \item 不是全行范围审批
    \end{itemize}

    \textbf{B选项错误。}
    这个说法不准确：
    \begin{itemize}
        \item 流动性期限由巴塞尔框架规定
        \item 银行不能自行决定
    \end{itemize}

    \textbf{C选项错误。}
    这个说法不准确：
    \begin{itemize}
        \item FRTB对交易账户和银行账户划分更严格
        \item 不是更灵活
    \end{itemize}

    \textbf{D选项正确。}
    这个说法正确：
    \begin{itemize}
        \item 即使获准使用内部模型法
        \item 也必须使用标准法计算最低资本要求
        \item 作为资本要求的底线
    \end{itemize}
\end{explanation}

\checklist{FRTB框架的主要特征}



\begin{question}
    A country that is largely dependent on commodity exports is experiencing a period of very weak economic conditions due to a significant drop in key commodity prices as well as an overall increase in unemployment. A banking supervisor is concerned about the impact of the current market conditions on the country's banks and is using early warning indicators (EWIs) to detect potential adverse changes in the liquidity risk of the banks. Which of the following observations about a bank would be the greatest cause for concern to the supervisor?
\end{question}

\begin{choices}
    \item The bank's capital reserves have increased above its countercyclical buffer
    \item The bank has shifted some of its funding sources from short-term repo agreements to longer-term liabilities
    \item The bank's net stable funding ratio has increased
    \item The bank's Basel II leverage ratio has decreased
\end{choices}

\begin{correctanswer}
    D
\end{correctanswer}

\begin{explanation}
    \textbf{A选项错误。}
    这个说法不准确：
    \begin{itemize}
        \item 资本储备增加是积极信号
        \item 表明银行抗风险能力增强
    \end{itemize}

    \textbf{B选项错误。}
    这个说法不准确：
    \begin{itemize}
        \item 转向长期负债是积极信号
        \item 降低短期融资风险
    \end{itemize}

    \textbf{C选项错误。}
    这个说法不准确：
    \begin{itemize}
        \item 净稳定资金比率增加是积极信号
        \item 表明流动性状况改善
    \end{itemize}

    \textbf{D选项正确。}
    这个说法正确：
    \begin{itemize}
        \item 杠杆率下降是严重警告信号
        \item 表明银行资本充足性恶化
        \item 在经济下行期尤其危险
    \end{itemize}
\end{explanation}

\checklist{银行早期预警指标分析}



\newpage





\end{document}